\documentclass[12pt]{article}
\usepackage[utf8]{inputenc}
\usepackage[T1]{fontenc}
\usepackage[margin=1in]{geometry}
\usepackage{setspace}
\usepackage{xcolor}
\usepackage[colorlinks=true, urlcolor=blue!60!black, linkcolor=blue!60!black]{hyperref}
\setlength{\parindent}{2em}
\setlength{\parskip}{1.25em}
\renewcommand{\baselinestretch}{1.0}
\begin{document}
\begin{center}
{\LARGE \textit{Cruzan v.\ Director, Missouri Department of Health}, 497 U.S.\ 261 (1990)\par}
\vspace{0.8em}
{\large Case Brief --- Four-Page Edition\par}
\end{center}
\vspace{1em}

\section{Memory Jogger}

The Court assumed that a competent person has a Fourteenth Amendment liberty interest in refusing lifesaving nutrition and hydration, but held that a State may require ``clear and convincing evidence'' of an incompetent patient's wishes before family members may withdraw life support (\href{https://supreme.justia.com/cases/federal/us/497/261/}{Justia}).

\section{Detailed Case Facts}

In January 1983, 25-year-old Nancy Beth Cruzan lost control of her car and was thrown face-down into a ditch. Paramedics restored her heartbeat and breathing, but after 12 to 14 minutes without oxygen she fell into a persistent vegetative state, with no significant cognitive function. Surgeons implanted a gastrostomy feeding tube; experts testified she could be kept alive for another 30 years in a Missouri state hospital, at the State's expense. Her parents and co-guardians, Lester and Joyce Cruzan, asked the hospital to end artificial nutrition and hydration. Because removal would cause her death, the hospital refused without court approval. \textit{Cruzan} was the Court's first ``right to die'' case, arriving fourteen years after \textit{In re Quinlan}, and it helped prompt the Patient Self-Determination Act of 1990.

\section{Procedural History}

The trial court authorized withdrawal, relying on a former housemate's testimony that Nancy had said she would not want to live unless she could live ``at least halfway normally.'' The Missouri Supreme Court reversed, 4--3, calling those statements ``unreliable'' and requiring living-will formalities or ``clear and convincing, inherently reliable evidence'' of the patient's wishes. The U.S.\ Supreme Court affirmed. (After Missouri withdrew, new testimony from coworkers satisfied the standard; the tube was removed and Nancy died on December 26, 1990.)

\section{Judicial Votes}

\textbf{5--4}, affirmed (June 25, 1990).
\begin{itemize}
\item \textbf{Majority:} Chief Justice Rehnquist, joined by Justices White, O'Connor, Scalia, and Kennedy.
\item \textbf{Concurrences:} Justice O'Connor; Justice Scalia.
\item \textbf{Dissents:} Justice Brennan, joined by Justices Marshall and Blackmun; and Justice Stevens, writing alone.
\end{itemize}

\section{Holding}

A competent person has a due process liberty interest in refusing unwanted medical treatment, and the Court assumed that it extends to lifesaving artificial nutrition and hydration. But the Constitution does not forbid a State from requiring clear and convincing evidence of an incompetent patient's wishes before a surrogate may withdraw treatment, nor does it require a State to accept the ``substituted judgment'' of close family members. Footnote 12 reserved whether a State must honor a surrogate the patient formally designated while competent.

\section{Analysis of Opinions}

\textbf{Majority (Rehnquist).} Surveying the common law of informed consent and precedents such as \textit{Jacobson v.\ Massachusetts}, the Court recognized a ``liberty interest'' (not a privacy right) to be balanced against state interests. An incompetent patient's choice must be exercised by a surrogate, and Missouri could demand heightened proof to serve its unqualified interest in preserving life without ``quality of life'' judgments, to guard against family abuse, and because surrogate hearings may not be adversarial. Under \textit{Addington} and \textit{Santosky}, proof standards allocate the risk of error, and an erroneous decision to withdraw treatment is irrevocable.

\textbf{O'Connor (concurring).} Tube feeding involves physical intrusion, so the Due Process Clause protects the decision to refuse it; and honoring a duly appointed health-care proxy ``may well be constitutionally required.''

\textbf{Scalia (concurring).} ``The federal courts have no business in this field.'' The common law always let the State prevent suicide, including by refusing sustenance; the action/inaction line is specious.

\textbf{Brennan, joined by Marshall and Blackmun (dissenting).} Freedom from unwanted treatment is a \textit{fundamental} right. The State's only legitimate interest was accuracy, and error is irrevocable in both directions: wrongly maintaining a vegetative patient degrades her dignity and prolongs her family's anguish.

\textbf{Stevens (dissenting).} ``Choices about death touch the core of liberty.'' Nancy's parents, the trial judge, and the guardian ad litem agreed withdrawal served her best interests; Missouri imposed a ``sectarian definition of life.''

\section{Examples: Future Applications}

\begin{enumerate}
\item \textbf{Same Side: Ventilator Rule.} A State bars withdrawal of a ventilator from a comatose patient absent an advance directive or clear and convincing proof. \textit{Cruzan} allows the State to allocate the risk of error toward life, even over unanimous family belief.
\item \textbf{Same Side: Assisted Suicide.} A ban on prescribing lethal medication to a competent, terminally ill adult was upheld in \textit{Washington v.\ Glucksberg} (1997), which distinguished refusing treatment from assistance in suicide.
\item \textbf{Opposite Side: Competent Refusal.} A state hospital force-feeds a competent cancer patient who repeatedly refuses. Five Justices recognized, and the majority assumed, her right to refuse.
\item \textbf{Opposite Side: Ignoring a Living Will.} A law requiring hospitals to ignore a clear written directive against tube feeding extinguishes the liberty interest without any evidentiary justification.
\item \textbf{Fence-Sitter: Designated Proxy.} A patient names her sister as health-care proxy but leaves no instructions on feeding tubes. Footnote 12 left the question open, and O'Connor suggested deference may be constitutionally required.
\end{enumerate}

\section{Critique}

\textbf{False neutrality.} Critics, echoing Brennan, argue that maintaining a permanently unconscious patient by surgically implanted tubes is an ongoing intervention, not a neutral ``status quo,'' and that invoking \textit{Addington} and \textit{Santosky}---cases protecting individuals against the State---inverted the protective function of proof standards.

\textbf{From both sides.} Originalists, following Scalia, say the Court assumed an unenumerated substantive right without grounding in text or history. Progressive and bioethical scholars say the Court downgraded bodily autonomy to a balanced ``liberty interest,'' favoring those sophisticated enough to draft living wills and leaving young accident victims like Nancy to indefinite maintenance.

\section{Key Quotations}

\begin{quote}
``The principle that a competent person has a constitutionally protected liberty interest in refusing unwanted medical treatment may be inferred from our prior decisions.''
\end{quote}
\noindent --- Chief Justice Rehnquist, for the Court, 497 U.S.\ at 278.

\begin{quote}
``An erroneous decision not to terminate results in a maintenance of the status quo; the possibility of subsequent developments such as advancements in medical science, the discovery of new evidence regarding the patient's intent, changes in the law, or simply the unexpected death of the patient despite the administration of life-sustaining treatment, at least create the potential that a wrong decision will eventually be corrected or its impact mitigated. An erroneous decision to withdraw life-sustaining treatment, however, is not susceptible of correction.''
\end{quote}
\noindent --- Chief Justice Rehnquist, for the Court, 497 U.S.\ at 283.

\begin{quote}
``Requiring a competent adult to endure such procedures against her will burdens the patient's liberty, dignity, and freedom to determine the course of her own treatment. Accordingly, the liberty guaranteed by the Due Process Clause must protect, if it protects anything, an individual's deeply personal decision to reject medical treatment, including the artificial delivery of food and water.''
\end{quote}
\noindent --- Justice O'Connor, concurring, 497 U.S.\ at 289.

\end{document}
